% ============================================================================
\section{From Language to Logic}
\label{sec:logic}
% Source: assignments/07_logic/index.html
% ============================================================================

Problems 1--5 are programming problems (formulas written with \texttt{logic.py}). We list the formulas in
mathematical notation so the report can be read on its own; the code and grader output are in the notebook.

% ----------------------------------------------------------------------------
\subsection{Propositional logic}
\codepart{7.1}
\begin{enumerate}[label=(1\alph*)]
  \item ``If it's summer and we're in California, then it doesn't rain.'' \quad \todo{formula}
  \item ``It's wet if and only if it is raining or the sprinklers are on.'' \quad \todo{formula}
  \item ``Either it's day or night (but not both).'' \quad \todo{formula}
\end{enumerate}

% ----------------------------------------------------------------------------
\subsection{First-order logic}
\codepart{7.2}
\begin{enumerate}[label=(2\alph*)]
  \item ``Every person has a parent.'' \quad \todo{formula}
  \item ``At least one person has no children.'' \quad \todo{formula}
  \item Define $\text{Father}(x,y)$ in terms of $\text{Male}$ and $\text{Parent}$. \quad \todo{formula}
  \item Define $\text{Granddaughter}(x,y)$ in terms of $\text{Female}$ and $\text{Child}$. \quad \todo{formula}
\end{enumerate}

% ----------------------------------------------------------------------------
\subsection{Liar puzzle}
\codepart{7.3a}
Suspects: John, Susan, Mark, Nicole. Mark: ``It wasn't me!'' John: ``It was Nicole!'' Nicole: ``It wasn't me, it
was Susan!'' Susan: ``Nicole's a liar.'' Exactly one person tells the truth, and exactly one person crashed the server.

\todo{the six formulas and the answer returned by \texttt{grader.py 3a-run}}

% ----------------------------------------------------------------------------
\subsection{Odd and even integers}
\codepart{7.4a}
\todo{the six axioms in mathematical notation and the result of \texttt{grader.py 4a-run}}

% ----------------------------------------------------------------------------
\subsection{Semantic parsing (extra credit)}
\codepart{7.5}
Grammar rules for: \emph{every \$Noun \$Verb some \$Noun}; \emph{there is some \$Noun that every \$Noun \$Verb};
\emph{if a \$Noun \$Verb a \$Noun then the former \$Verb the latter}.

% ----------------------------------------------------------------------------
\subsection{Explainability of logic-based systems}

The four considerations: \emph{respect}; \emph{assessing the fairness of rules}; \emph{contesting and correcting
decisions}; \emph{ability to change user behaviour}. Context: the liar puzzle of Problem~3, where a conclusion could
have real consequences (e.g.\ someone being fired).

\problempart{6a}{What counts as an explanation}
\question{Besides the final answer, what does a general first-order logic system, like the one used in Problem 3,
return that could count as an explanation?}
\expect{1--2 sentences.}
\begin{solution}
\todo{answer}
\end{solution}

\problempart{6b}{Why it may not be adequate}
\question{Why might this explanation not be adequate? (Hint: model misspecification, or how the question is framed.)}
\expect{1--2 sentences.}
\begin{solution}
\todo{answer}
\end{solution}

\problempart{6c}{Which considerations are met}
\question{Are any of the four explainability considerations met by the explanations a logic-based system produces?}
\expect{2--4 sentences.}
\begin{solution}
\todo{answer}
\end{solution}

% ----------------------------------------------------------------------------
\subsection{Applications of soundness and completeness}

Sound: $\{f: KB\vdash f\}\subseteq\{f: KB\vDash f\}$. \quad Complete: $\{f: KB\vdash f\}\supseteq\{f: KB\vDash f\}$.

\problempart{7a}{Is an LLM sound or complete?}
\question{Treat whatever GPT-3.5 generates as what it ``derives''. Given that LLMs hallucinate, is GPT-3.5 sound,
complete, both or neither? Use examples and the definitions.}
\expect{2--3 sentences for each property (4--6 in total).}
\begin{solution}
\todo{answer}
\end{solution}

\problempart{7b}{Safety-critical deployment}
\question{System A is sound but not complete, and System B is complete but not sound. Which would you deploy in a
safety-critical setting, where acting on a false conclusion has severe consequences, and why?}
\expect{2--3 sentences tied to the definitions.}
\begin{solution}
\todo{answer}
\end{solution}

\problempart{7c}{With a separate verifier}
\question{With a separate verifier available (e.g.\ a human in the loop) and the goal of capturing as much true
information as possible, which system would you choose, and why?}
\expect{2--3 sentences tied to the definitions.}
\begin{solution}
\todo{answer}
\end{solution}
